\section{When the subsidiary earns the internal margin}
\label{sec:worked-intragroup-upstream}

Now let the subsidiary sell the inventory to the parent. The group must still
eliminate the full internal sale, but the profit being removed was earned in
the company partly owned by NCI. Applying the minority percentage to that
company's unadjusted legal profit would allocate profit that the group has not
yet earned from an outside customer.

Keep the acquisition on 1 July 2025 unchanged. On 1 September, after control
has begun, the subsidiary buys inventory for 20 cash from an outside vendor.
On 1 October it sells the lot to the parent for 30 on credit. Buying the lot
after control matters: inventory is not inserted into the acquired net assets,
and the original goodwill and acquisition NCI remain unchanged. The subsidiary's
cash falls by 20. This external payment survives consolidation even if the
parent later settles the internal payable of 30. Settlement creates a payment
by the parent and a receipt by the subsidiary; both internal cash flows are
eliminated in full.

Before any outside resale, the parent's legal inventory is 30 and the
subsidiary has legal trading profit 10. Consolidation removes the reciprocal
receivable and payable, eliminates internal revenue 30, and reduces inventory
by the unrealised margin 10. Group inventory is therefore 20. Neither the
inventory adjustment nor the reciprocal-balance elimination is multiplied by
the parent's 75\% ownership. Full elimination precedes the allocation of the
subsidiary's adjusted earnings between parent owners and NCI
\citep[paras.~B86, B89--B94]{aasb10control2026}.

\subsection{Tax recognition and attribution answer different questions}

Continue with explicitly supplied separate tax bases and 25\% rates. The
subsidiary owes current tax $5/2$ on its legal margin. If the parent's inventory
tax base remains 30, the difference from the group carrying amount 20 can
support a deferred tax asset of $5/2$ when the recognition conditions are met.
The comparison assumes those conditions in one branch and explicitly withholds
recognition in another. The tax base and recognition conditions remain distinct
from the inventory elimination \citep[paras.~11, 24 and illustrative
example~B11]{aasb112dated}.

For this example, attribute the inventory elimination and its deferred tax
effect to the subsidiary seller's earnings. This is a supplied attribution
premise; the cited paragraphs do not themselves determine the example's tax
sharing. With recognised deferred tax, removing profit 10 and adding tax income
$5/2$ offsets the subsidiary's after-tax legal margin $15/2$. Its consolidated
profit returns to $63/2$, the acquired-basis operating result already derived.
NCI receives one quarter, or $63/8$, and parent owners receive $189/8$.

The group now holds cash 250, inventory 20, equipment 81, goodwill $135/4$
and the inventory deferred tax asset $5/2$. Total assets are $1549/4$;
liabilities are 27, including the unpaid current tax. Equity is $1441/4$.
If no inventory deferred tax asset is recognised, the current tax remains an
expense. Group profit becomes 29, of which NCI receives $29/4$ and parent
owners $87/4$. Assigning that entire tax expense to parent owners would be
wrong under the stated subsidiary-attribution premise. If tax is unresolved,
the reference retains the pretax worksheet and the outstanding tax questions
but produces no final financial report.

\subsection{Outside resale releases the subsidiary's margin}

Let the parent sell a fraction $r$ of the lot outside the group on 1 December
for $40r$ cash, where this finite comparison permits $r=0$, $1/2$ or $1$.
Its legal cost is $30r$ and its current tax is $5r/2$, so its own after-tax
trading contribution is $15r/2$. The subsidiary's original margin 10 is now
realised only in proportion $r$. Write $d$ for the remaining recognised
inventory deferred tax asset: $d=5(1-r)/2$ in the recognised branch and zero
in the unrecognised branch. Under the supplied attribution, the two trading
contributions and resulting NCI profit are
\[
 \pi_S^{\mathrm{trade}}=10r-\frac52+d,\qquad
 \pi_P^{\mathrm{trade}}=\frac{15r}{2},\qquad
 \pi_{\mathrm{NCI}}=\frac14
       \left(\frac{63}{2}+\pi_S^{\mathrm{trade}}\right).
\]
The subsidiary contribution combines realised margin, its unchanged current
tax, and the deferred benefit still recognised for unsold inventory. With
recognised deferred tax it simplifies to $15r/2$. Half resale therefore adds
$15/4$ of profit to each company: group profit is 39 and NCI profit is $141/16$.
The group holds cash 270 and inventory 10; total assets are 396, liabilities
$113/4$ and equity $1471/4$. The external vendor payment and outside customer
receipt remain visible separately in the cash record.

Equal recognised-tax trading contributions happen to make several profit
totals match the downstream example. That numerical coincidence is not a
reason to forget which company earned the margin. The unrecognised-tax branch,
the separate entity accounts and the carried opening equity distinguish the
two directions. Complete legal journals, account worksheets, cash origins and
attribution results are compared across the finite cases.\footnote{The exact
upstream comparisons, malformed-request controls and protected sources are in
\path{data/accounting-master/acquisition-intragroup-upstream-20261027/}.
The dates, lot, customer-control facts, rights and tax premises are synthetic;
no actual source or accounting admission is inferred.}
The remaining question is how the subsidiary-attributed margin and NCI move
through retained profit when the outside sale occurs in a later year.
